%! Exponential Growth and Decay — Unit 7, Lesson 7.1 — Guided Notes (Answer Key)
\documentclass[10pt]{article}
\usepackage{algebra2-article}
\usepackage{algebra2-key}   % pulls in -boxes; do NOT also load -boxes
\newcommand{\TallMath}[1]{$\displaystyle #1\rule[-1.4em]{0pt}{3.2em}$}
\newcommand{\lgrid}[4]{%
  \draw[step=1, linegray!40, very thin] (#1,#3) grid (#2,#4);
  \draw[->, semithick, charcoal] (#1,0)--(#2,0) node[below right, font=\scriptsize]{$x$};
  \draw[->, semithick, charcoal] (0,#3)--(0,#4) node[left, font=\scriptsize]{$y$};}
% Exponentials are plotted as a*e^{(ln b)x}: pgfmath's pow() is unreliable for negative and
% fractional exponents, but exp() is not. #2 is ln(b).
% ln2 = 0.693147   ln3 = 1.098612   ln4 = 1.386294
\newcommand{\expcurve}[4]{\draw[very thick, forest, domain=#3:#4, samples=120, smooth]
  plot (\x,{#1*exp((#2)*(\x))});}
% Vocab answer for the key: bold forest term, then the filled definition on a write-line.
% The leading and trailing \par are required: \ansline ends with \dotfill but does NOT end the
% paragraph, so without them each term label gets pulled onto the previous answer's line.
\newcommand{\vocabans}[2]{%
  \par\noindent\textbf{\textcolor{forest}{#1:}}\\[1pt]\ansline{#2}\par}

\begin{document}

\pageheader{Unit 7, Lesson 7.1}{Guided Notes --- Answer Key}
\namedateperiod

\vspace{0.08in}

\begin{objectivebox}
\textbf{By the end of this lesson, I will be able to\dots}
\begin{itemize}\setlength{\itemsep}{2pt}
  \item \emph{build} the equation $y=ab^{x}$ from a table, from two points, and from a written description;
  \item translate a \emph{percent} into a \emph{factor} --- ``up $8\%$'' means $b=1.08$ and ``down $15\%$'' means $b=0.85$ --- and explain why the percent itself is never the base;
  \item evaluate a model I built and interpret the answer in the context it came from.
\end{itemize}
\end{objectivebox}

\vspace{0.05in}

\begin{vocabbox}
\small
Fill in each term as we build it together.
\par\vspace{2pt}   % \par is required: \vocabans's \noindent is a no-op mid-paragraph
\vocabans{Initial value $a$}{the amount the model starts with --- the output at $x=0$, and the $y$-intercept of the graph}
\vocabans{Growth (or decay) factor $b$}{the number you multiply by for each one-unit step; $b>1$ grows, $0<b<1$ decays}
\vocabans{Growth (or decay) rate $r$}{the percent change per step, written as a decimal --- $8\%$ becomes $r=0.08$}
\vocabans{Exponential growth model $y=a(1+r)^{t}$}{a quantity that increases by the same \emph{percent} each period; the factor is $b=1+r$}
\vocabans{Exponential decay model $y=a(1-r)^{t}$}{a quantity that decreases by the same \emph{percent} each period; the factor is $b=1-r$, the fraction that \emph{remains}}
\end{vocabbox}

\begin{hookbox}
Two towns publish their plans for the next ten years.
\begin{center}
\small
\renewcommand{\arraystretch}{1.25}
\begin{tabularx}{0.92\linewidth}{@{}l X@{}}
\hline
\textbf{Riverton} & Population $8{,}000$ today. \emph{``We expect to gain $300$ people every year.''} \\
\textbf{Kellsboro} & Population $8{,}000$ today. \emph{``We expect to grow by $4\%$ every year.''} \\
\hline
\end{tabularx}
\end{center}
\begin{itemize}\setlength{\itemsep}{1pt}
  \item Which town's plan is a \textbf{linear} model, and which is \textbf{exponential}? How can you tell without doing any arithmetic? \ansline{Riverton is linear --- it \emph{adds} the same number of people every year (constant difference). Kellsboro is exponential --- it \emph{multiplies} by the same factor every year (constant ratio).}
  \item Kellsboro grows $4\%$ this year. Next year it grows $4\%$ again --- of \emph{what}? \ansline{Of the \emph{new}, larger population --- not of the original $8{,}000$. That is exactly what ``multiply by the same factor each time'' means.}
\end{itemize}
Yesterday somebody handed you $y=ab^{x}$ and you read it. Today you have to \textbf{write} it. And the
whole thing comes down to one question: \emph{when a sentence says $4\%$, what number goes in the base?}
\end{hookbox}

\begin{notesbox}{1. You Only Ever Need Two Numbers --- and a Table Hands You Both}
An exponential model is $y=ab^{x}$, so building one means finding exactly two numbers:
\begin{center}
\renewcommand{\arraystretch}{1.35}
\begin{tabularx}{0.94\linewidth}{@{}l l X@{}}
\hline
$a$ & the \textbf{initial value} & the output when $x=$ \ans{0} --- where the model \emph{starts} \\
$b$ & the \textbf{factor} & what you \textbf{multiply} by each time $x$ goes up by $1$ \\
\hline
\end{tabularx}
\end{center}

\vspace{3pt}
\textbf{From a table} (with $x$-values one step apart), both numbers are sitting right there:
\begin{center}
\renewcommand{\arraystretch}{1.3}
\begin{tabular}{|c|c|c|c|c|c|}
\hline
$x$ & $0$ & $1$ & $2$ & $3$ & $4$ \\ \hline
$y$ & $3$ & $12$ & $48$ & $192$ & $768$ \\ \hline
\end{tabular}
\end{center}
\begin{enumerate}[leftmargin=1.7em, itemsep=3pt]
  \item \textbf{Read $a$} straight off the $x=0$ column: $a=$ \ans{3}.
  \item \textbf{Divide to get $b$:} $\frac{12}{3}=$ \ans{4}, \; $\frac{48}{12}=$
        \ans{4}, \; $\frac{192}{48}=$ \ans{4}. \; Same every time, so the table
        \emph{is} exponential and $b=$ \ans{4}.
  \item \textbf{Write it:} $y=$ \ans{$3\cdot 4^{x}$}
  \item \textbf{Check} with a row you did not use: at $x=4$ the model gives \ans{$3\cdot 4^{4}=768$},
        and the table says \ans{768}. \; Match? \ans{yes}
\end{enumerate}

\vspace{3pt}
\textbf{Your turn.} Build the model for this table.
\begin{center}
\renewcommand{\arraystretch}{1.3}
\begin{tabular}{|c|c|c|c|c|}
\hline
$x$ & $0$ & $1$ & $2$ & $3$ \\ \hline
$y$ & $800$ & $200$ & $50$ & $12.5$ \\ \hline
\end{tabular}
\end{center}
$a=$ \ans{800} \quad $b=$ \ans{$\frac14$ (or $0.25$)} \quad Growth or decay? \ans{decay} \quad
$y=$ \ans{$800\left(\frac14\right)^{x}$}

\vspace{2pt}
\emph{Careful:} always divide \textbf{a later output by an earlier one}. Doing it upside down here
gives $4$ and turns a shrinking table into a growing model.
\end{notesbox}

\begin{notesbox}{2. Percent $\to$ Factor: the Translation This Whole Lesson Runs On}
A sentence almost never hands you $b$. It hands you a \textbf{percent rate} $r$, and you have to
convert. Go back to the Warm-Up: a \$$50$ jacket that goes \textbf{up} $8\%$ is now worth $108\%$ of
\$$50$, so you multiply by \ans{1.08}. A \$$50$ jacket that goes \textbf{down} $15\%$ is now worth
\ans{85}\% of \$$50$, so you multiply by \ans{0.85}.
\[
  \textbf{Growth:}\quad b = 1 + r
  \qquad\qquad
  \textbf{Decay:}\quad b = 1 - r
\]
\begin{center}
\renewcommand{\arraystretch}{1.45}
\begin{tabularx}{\linewidth}{@{}X l l X@{}}
\hline
\textbf{The sentence says\dots} & \textbf{Rate $r$} & \textbf{Grows or shrinks?} & \textbf{Factor $b$} \\
\hline
increases by $8\%$ each year & \ans{0.08} & \ans{grows} & \ans{$1.08$} \\
decreases by $15\%$ each year & \ans{0.15} & \ans{shrinks} & \ans{$0.85$} \\
increases by $3.5\%$ each year & \ans{0.035} & \ans{grows} & \ans{$1.035$} \\
loses $40\%$ of its value each hour & \ans{0.40} & \ans{shrinks} & \ans{$0.60$} \\
$100\%$ of it remains each day & \ans{0} & \ans{neither} & \ans{$1$ --- illegal!} \\
\textbf{doubles} each year & \ans{$1$ ($100\%$)} & \ans{grows} & \ans{$2$} \\
\hline
\end{tabularx}
\end{center}

\vspace{2pt}
\begin{tcolorbox}[colback=redbg, colframe=redacc, boxrule=0.8pt, arc=1.5mm,
    left=2.5mm, right=2.5mm, top=1.5mm, bottom=1.5mm]
\textbf{The two errors everybody makes today.} Both are ``I used the percent as the base.''
\begin{itemize}[leftmargin=1.4em, itemsep=2pt]
  \item Writing $b=0.08$ for ``up $8\%$.'' That base is \emph{less} than $1$, so it is a
        \ans{decay} model --- and it says only $8\%$ of the amount survives, a
        \ans{92}\% \emph{loss}.
  \item Writing $b=1.15$ for ``down $15\%$.'' That base is \emph{greater} than $1$, so it makes the
        quantity \ans{grow} --- it models a $15\%$ \ans{increase} instead.
\end{itemize}
\textbf{The two-second check:} after you write $b$, ask \emph{is $b$ on the correct side of $1$?}
Growth needs $b>1$; decay needs $0<b<1$.
\end{tcolorbox}
\end{notesbox}

\begin{notesbox}{3. From a Story to a Model --- and Back to the Story}
Same three moves every time: pull out $a$, translate the percent into $b$, then write $y=ab^{x}$.

\vspace{4pt}
\textbf{Worked example --- growth.} A town of $12{,}000$ people grows by $3\%$ each year. Write a
model for the population $P$ after $t$ years, then predict the population in $10$ years.
\begin{enumerate}[leftmargin=1.7em, itemsep=3pt]
  \item Initial value: $a=$ \ans{$12{,}000$} \quad (the population at $t=0$, \emph{today}).
  \item Rate $r=$ \ans{0.03}, and the town is \ans{growing}, so
        $b=1+$ \ans{0.03} $=$ \ans{1.03}.
  \item Model: $P(t)=$ \ans{$12{,}000(1.03)^{t}$}
  \item Predict: $P(10)=$ \ans{$12{,}000(1.03)^{10}=12{,}000(1.3439)$} $\approx$ \ans{$16{,}127$} people.
  \item Interpret in a sentence: \ansline{If the $3\%$ yearly growth keeps up, the town will have about $16{,}127$ people ten years from now --- roughly $4{,}100$ more than today.}
\end{enumerate}

\vspace{3pt}
\textbf{Worked example --- decay.} A car is bought for \$$24{,}000$ and loses $12\%$ of its value
every year. Write a model for the value $V$ after $t$ years, then find its value after $5$ years.
\begin{enumerate}[leftmargin=1.7em, itemsep=3pt]
  \item $a=$ \ans{$24{,}000$} \quad $r=$ \ans{0.12} \quad
        $b=1-$ \ans{0.12} $=$ \ans{0.88}
  \item Model: $V(t)=$ \ans{$24{,}000(0.88)^{t}$}
  \item $V(5)=$ \ans{$24{,}000(0.88)^{5}=24{,}000(0.5277)$} $\approx$ \ans{\$$12{,}666$}
  \item Interpret in a sentence: \ansline{Five years after it was bought, the car is worth about \$$12{,}666$ --- a little more than half its original price.}
  \item Why is $b=0.88$ and \emph{not} $0.12$? \ansline{$0.12$ is the fraction the car \emph{loses}. The base has to be the fraction that \emph{remains}: $100\%-12\%=88\%$, so $b=0.88$.}
\end{enumerate}

\vspace{2pt}
\textbf{The check-your-model routine.} Before you use a model you just wrote, run all three:
\begin{enumerate}[label=(\arabic*), leftmargin=1.9em, itemsep=2pt]
  \item Does $y$ at $x=0$ give back the \ans{starting} amount from the story?
  \item Is $b$ on the correct \ans{side} of $1$ for growth vs.\ decay?
  \item Does one step forward change the answer by the right \ans{percent}?
\end{enumerate}
\end{notesbox}

\begin{notesbox}{4. Building From Two Points --- Dividing Makes $a$ Disappear}
Sometimes all you get is two points on the curve.

\vspace{3pt}
\textbf{Case 1: one of them is the $y$-intercept.} The curve passes through $(0,5)$ and $(3,40)$.
\begin{itemize}[leftmargin=1.5em, itemsep=3pt]
  \item The point $(0,5)$ gives $a$ immediately: $a=$ \ans{5}.
  \item Substitute the other point into $y=ab^{x}$: \; $40 = 5\,b^{\,3}$, so $b^{3}=$
        \ans{8} and $b=$ \ans{2}. \;
        \emph{(A cube root --- Unit 6 machinery, nothing new.)}
  \item Model: $y=$ \ans{$5\cdot 2^{x}$}
\end{itemize}

\vspace{4pt}
\textbf{Case 2: neither point is the $y$-intercept.} Read the two labeled points off the graph.
\begin{center}
\begin{minipage}[c]{0.46\linewidth}
\centering
\begin{tikzpicture}[scale=0.44]
  \lgrid{-4}{4}{-1}{9}
  \begin{scope}
    \clip (-4,-1) rectangle (4,9);
    \expcurve{2}{0.693147}{-4}{2.17}
  \end{scope}
  \draw[navy, dashed, semithick] (-4,0)--(4,0);
  \fill[charcoal] (-1,1) circle (3pt) node[below right, font=\scriptsize]{$(-1,1)$};
  \fill[charcoal] (2,8) circle (3pt) node[left, font=\scriptsize]{$(2,8)$};
\end{tikzpicture}
\end{minipage}
\begin{minipage}[c]{0.50\linewidth}
\small
\textbf{Divide the two equations.} Writing both points into $y=ab^{x}$ gives
$1=a\,b^{-1}$ and $8=a\,b^{2}$. Dividing the second by the first cancels $a$:
\[
  \frac{8}{1} \;=\; \frac{a\,b^{2}}{a\,b^{-1}} \;=\; b^{\,3}
\]
so $b^{3}=$ \ans{8} and $b=$ \ans{2}.
\end{minipage}
\end{center}
\begin{itemize}[leftmargin=1.5em, itemsep=3pt]
  \item Now put $b$ back into \emph{either} point to get $a$: \; $8=a\cdot$ \ans{4}, so
        $a=$ \ans{2}.
  \item Model: $y=$ \ans{$2\cdot 2^{x}$}. \; Check it against the graph: the curve should cross the
        $y$-axis at \ans{$(0,2)$}, and it does.
  \item \textbf{Why the exponent came out $3$:} the two $x$-values are \ans{three} steps apart,
        and each step multiplies by $b$ once.
\end{itemize}
\end{notesbox}

\begin{practicebox}
Build the model in each case. Write your answer in the form $y=ab^{x}$.
\begin{enumerate}[leftmargin=1.7em, itemsep=6pt]
  \item \textbf{From a table.}
    \begin{center}
    \renewcommand{\arraystretch}{1.25}
    \begin{tabular}{|c|c|c|c|c|}
    \hline
    $x$ & $0$ & $1$ & $2$ & $3$ \\ \hline
    $y$ & $4$ & $10$ & $25$ & $62.5$ \\ \hline
    \end{tabular}
    \end{center}
    $a=$ \ans{4} \quad $b=$ \ans{2.5} \quad $y=$ \ans{$4(2.5)^{x}$}

  \item \textbf{From a story.} A city currently uses $900$ million gallons of water a year and has
        cut its use by $4\%$ each year. \\[3pt]
        $W(t)=$ \ans{$900(0.96)^{t}$} \quad $W(5)=$ \ans{$\approx 733.8$} million gallons \\[3pt]
        What does $W(5)$ mean? \ansline{If the $4\%$ yearly cut continues, in five years the city will use about $733.8$ million gallons --- roughly $166$ million gallons less than today.}

  \item \textbf{From two points.} A curve $y=ab^{x}$ passes through $(0,7)$ and $(2,63)$. \\[3pt]
        $a=$ \ans{7} \quad $b^{2}=$ \ans{9} \quad $b=$ \ans{3} \quad
        $y=$ \ans{$7\cdot 3^{x}$}

  \item A classmate models ``a savings account of \$$400$ that earns $6\%$ per year'' as
        $S(t)=400(0.06)^{t}$. Without computing anything, give \emph{one} reason this cannot be
        right. \ansline{The base $0.06$ is less than $1$, so the model \emph{shrinks} --- but the account is supposed to grow. (It actually models losing $94\%$ a year.) The base should be $1+0.06=1.06$.}
\end{enumerate}
\end{practicebox}

\begin{teachernote}
\textbf{The one sentence to protect:} \emph{the base is the fraction you end up with, not the
percent that changed.} Every error today is a violation of it --- $b=0.08$ for ``up $8\%$,''
$b=1.15$ for ``down $15\%$,'' and $b=0.12$ for ``loses $12\%$'' are all the same mistake. The
Warm-Up's ``$108\%$ of the old price'' line is the fix; keep pointing back at it.

\textbf{Section 2, row 5} (``$100\%$ of it remains'') is deliberately the illegal base $b=1$. It pays
off Lesson 7.0's restriction $b\ne1$ with a reason students can feel: nothing changes, so the graph
is a horizontal line and the function is not exponential at all. Spend thirty seconds, not five
minutes.

\textbf{Section 3, decay item 5} (``why $0.88$ and not $0.12$?'') is the item to cold-call. If the
class can answer it in their own words, the rest of the lesson is mechanical. Homework problem 5 and
the Activity's Tier A justification both re-ask it, so you will get two more reads on it.

\textbf{Section 4 is the hardest page.} Case 2's division step is the new idea --- $a$ cancels
because it is a common factor --- and the exponent $3$ comes from the \emph{gap} between the
$x$-values, which is the part students drop. If time is short, do Case 1 fully and walk Case 2
aloud with the graph; Tier E of the Activity is where it gets real practice, and Lesson 7.5's
regression is where it gets automated.

\textbf{Rounding convention:} money to the nearest cent or dollar, populations to whole people, and
say so out loud --- students who report $16{,}126.997$ people should be redirected, not marked wrong.
\end{teachernote}

\end{document}
